% chapters/requirements/rf_annotations.tex
% RF-10: Anotaciones
% ============================================================================

\item \textbf{Anotaciones.} \label{rf:annotations}
Las anotaciones son los registros que los correctores crean sobre las instancias de examen
durante la corrección. Existen tres tipos: texto enriquecido (almacenado como
\textit{payload} opaco que el \dfn{backend} no interpreta), trazos de \textit{stylus}
(almacenados como \acs{json} de trazos con coordenadas, presión y color) y audio (con un
límite de duración configurable por organización). Una anotación puede estar asociada a un
problema concreto o ser general de la instancia. Los \textit{payloads} binarios se
almacenan en MinIO; los metadatos en la base de datos.

\begin{functional}

    % RF-10.1 ---------------------------------------------------------------
    \item \textbf{Creación de anotación.} \label{rf:annotations:create}
    El sistema debe permitir a un corrector crear una anotación sobre una instancia
    asignada.

    \textbf{Entrada:} petición POST al recurso de anotaciones de la instancia con los datos
    de la anotación: tipo (texto, \textit{stylus} o audio), \textit{payload} (contenido
    textual, \acs{json} de trazos o fichero de audio), ubicación en el \acs{pdf} (página y
    coordenadas) y, opcionalmente, el identificador del problema asociado.

    \textbf{Proceso:} el sistema valida que el corrector tiene la instancia asignada (o el
    problema, si solo tiene asignados problemas concretos, vía \texttt{AssignmentRule}).
    Para anotaciones de audio, se valida que la duración no excede el límite configurado
    (por defecto, 60~minutos). El \textit{payload} se almacena en MinIO (en el
    \dfn{bucket} de la organización). Se registran el autor, la fecha y la ubicación de la
    anotación.

    \textbf{Salida:} respuesta con los datos de la anotación creada y código
    \acs{http}~201.

    % RF-10.2 ---------------------------------------------------------------
    \item \textbf{Modificación de anotación.} \label{rf:annotations:update}
    El sistema debe permitir a un corrector modificar una anotación existente,
    sobreescribiendo su contenido.

    \textbf{Entrada:} petición PUT al recurso de la anotación con el \textit{payload}
    actualizado.

    \textbf{Proceso:} el sistema sobreescribe el contenido anterior de la anotación. No se
    mantiene historial de versiones; el estado anterior se pierde. Si la anotación es de
    audio y el nuevo \textit{payload} también, se reemplaza el fichero en MinIO. Solo el
    autor de la anotación puede modificarla.

    \textbf{Salida:} respuesta con los datos actualizados de la anotación.

    % RF-10.3 ---------------------------------------------------------------
    \item \textbf{Eliminación de anotación.} \label{rf:annotations:delete}
    El sistema debe permitir a un corrector eliminar una anotación que haya creado.

    \textbf{Entrada:} petición DELETE al recurso de la anotación.

    \textbf{Proceso:} el sistema elimina la anotación de la base de datos y el fichero
    asociado del almacenamiento de objetos.

    \textbf{Salida:} respuesta con código \acs{http}~204.

    % RF-10.4 ---------------------------------------------------------------
    \item \textbf{Consulta de anotaciones de una instancia.}
    \label{rf:annotations:list}
    El sistema debe permitir consultar las anotaciones de una instancia, con filtros por
    problema y página.

    \textbf{Entrada:} petición GET al recurso de anotaciones de la instancia con parámetros
    opcionales de filtrado: problema, página, tipo de anotación y autor.

    \textbf{Proceso:} el sistema recupera las anotaciones que cumplen los filtros. Los
    correctores acceden a las anotaciones de las instancias que tienen asignadas. Los
    estudiantes acceden a las anotaciones de su propia instancia solo durante la revisión y
    si el examen tiene activado el permiso de ver anotaciones.

    \textbf{Salida:} respuesta con la lista de anotaciones, cada una con sus metadatos y
    \textit{payload} (o referencia al \textit{payload} binario).

    % RF-10.5 ---------------------------------------------------------------
    \item \textbf{Anotación de calificación.} \label{rf:annotations:grade_mark}
    El sistema debe soportar una anotación especial con un indicador que la identifica como
    anotación de calificación, desencadenando la asignación automática de nota.

    \textbf{Entrada:} petición POST al recurso de anotaciones de calificación, con el
    \textit{payload} de la anotación y el identificador del problema.

    \textbf{Proceso:} se ejecuta \acs{ocr} sobre el \textit{payload} (trazos o texto) para
    extraer el valor numérico (véase \cref{rf:ingestion:ocr_grade}). La anotación se
    almacena como cualquier otra, pero con el indicador de calificación activo.

    \textbf{Salida:} respuesta con la anotación creada y la calificación asignada (o el
    resultado del \acs{ocr} si se requiere confirmación).

    % RF-10.6 ---------------------------------------------------------------
    \item \textbf{Almacenamiento de \textit{payloads} en MinIO.}
    \label{rf:annotations:storage}
    El sistema debe almacenar los \textit{payloads} de las anotaciones en el servicio de
    almacenamiento de objetos.

    \textbf{Entrada:} evento interno disparado por la creación de una anotación.

    \textbf{Proceso:} el fichero se almacena en MinIO en un \dfn{bucket} dedicado a
    anotaciones (dentro del \textit{namespace} de la organización), con una clave que
    incluye el identificador de la instancia y la anotación. Para los audios, se valida el
    formato y la duración antes del almacenamiento.

    \textbf{Salida:} confirmación del almacenamiento con la referencia interna al objeto en
    MinIO.

    % RF-10.7 ---------------------------------------------------------------
    \item \textbf{Asociación opcional a problema.} \label{rf:annotations:problem}
    El sistema debe permitir que una anotación esté asociada a un problema concreto o sea
    general de la instancia.

    \textbf{Entrada:} campo opcional de identificador de problema en la petición de creación
    o modificación de la anotación.

    \textbf{Proceso:} si se indica un problema, el sistema valida que el problema existe en
    el modelo de la instancia y que el corrector tiene asignado ese problema. Si no se
    indica problema, la anotación se considera general de la instancia (por ejemplo,
    comentarios en la portada o anotaciones transversales).

    \textbf{Salida:} anotación almacenada con o sin problema asociado.

\end{functional}
